\section{Local identities for every derivative of a spectral product}
\label{sec:spectral}

There is a useful second application of the distinction between a
multivariate germ and its one-dimensional restrictions. For polynomially
weighted products of shifted linear factors, every fixed spectral
derivative is a polynomial in the \emph{normalized factor exponents}.
Its highest-degree component is local and explicitly computable. This
produces finite identities between higher zeta derivatives even when the
individual values do not have short closed forms.

The first-derivative multiplicative anomaly is classical in substance;
Dowker treats linear factors and polynomial multiplicities
\cite{dowker}, and G\"urel gives general regularized-product formulas
\cite{gurel}, extending the Lerch--Mizuno setting \cite{mizuno}.
We retain that attribution. The all-order germ, the normalized-weight
polynomial formulation, and its explicit finite-difference consequences
are the proved extension developed here; the literature review does not
establish a priority claim for them.

\subsection{An exact multivariate meromorphic germ}

Let $P$ be a nonzero polynomial of degree $D$, let
$a_1,\ldots,a_m>0$, and initially define
\begin{equation}\label{sp:def}
 \cZ_P(\boldsymbol u;\boldsymbol a)
  =\sum_{n\ge0}P(n)\prod_{j=1}^m(n+a_j)^{-u_j},
 \qquad U=\sum_{j=1}^m u_j.
\end{equation}
Absolute convergence requires $\Rea U>D+1$; no separate positivity of
the $u_j$ is needed. Choose $b>0$ such that
$\max_j|a_j-b|<b$, for example $b=\max_j a_j$. Write
\begin{align}
 d_j&=a_j-b,& p_k&=\frac{P^{(k)}(-b)}{k!},&
 P(n)&=\sum_{k=0}^D p_k(n+b)^k,\label{sp:base}\\
 L_j(z)&=\log(1+d_jz),&
 L_{\boldsymbol u}(z)&=\sum_j u_jL_j(z),&
 C_\ell(\boldsymbol u)&=[z^\ell]e^{-L_{\boldsymbol u}(z)}.
 \label{sp:C}
\end{align}
The logarithms are their analytic germs at $z=0$. Explicitly,
\begin{equation}\label{sp:Cpoch}
 C_\ell(\boldsymbol u)=
 \sum_{q_1+\cdots+q_m=\ell}
       \prod_{j=1}^m\frac{(-1)^{q_j}d_j^{q_j}(u_j)_{q_j}}{q_j!}.
\end{equation}
For $\ell\ge1$ this is a polynomial of degree at most $\ell$ and has
zero constant term.

\begin{theorem}[Finite polar numerator, normally convergent regular part]
\label{sp:thm:germ}
The continuation of \eqref{sp:def} is
\begin{equation}\label{sp:continuation}
 \cZ_P(\boldsymbol u;\boldsymbol a)
   =\sum_{k=0}^D p_k\sum_{\ell\ge0}
          C_\ell(\boldsymbol u)\zeta(U+\ell-k,b).
\end{equation}
Near $\boldsymbol u=0$ it has the exact form
\begin{equation}\label{sp:germ}
 \boxed{\cZ_P(\boldsymbol u;\boldsymbol a)
             =\frac{\cR_P(\boldsymbol u)}U+\cH_P(\boldsymbol u),
 \qquad \cR_P=\sum_{k=0}^D p_k C_{k+1}.}
\end{equation}
Here $\cR_P(0)=0$, $\deg\cR_P\le D+1$, and $\cH_P$ is holomorphic.
An explicit regular part is
\begin{align}\label{sp:H}
 \cH_P(\boldsymbol u)=\sum_{k=0}^D p_k\Bigg\{&\zeta(U-k,b)
   +\sum_{\substack{\ell\ge1\\\ell\ne k+1}}
            C_\ell(\boldsymbol u)\zeta(U+\ell-k,b)\notag\\
   &+C_{k+1}(\boldsymbol u)g_b(U)\Bigg\}.
\end{align}
Every infinite sum displayed here converges normally, with every fixed
spectral derivative, away from its explicitly displayed poles.
\end{theorem}

\begin{proof}
Expand $(n+a_j)^{-u_j}$ as $(n+b)^{-u_j}$ times the binomial series
in $d_j/(n+b)$, then multiply the expansions. Choose a radius $\rho$
with $1/b<\rho<1/\max|d_j|$; when all $d_j=0$, any $\rho>1/b$
will do. On compact $\boldsymbol u$-sets the exponential in
\eqref{sp:C} is bounded on $|z|=\rho$, so
$|C_\ell(\boldsymbol u)|\le M\rho^{-\ell}$.
For sufficiently large $\ell$, the Dirichlet series for
$\zeta(U+\ell-k,b)$ and its fixed derivatives are bounded by a
constant times $b^{-\ell}$, allowing harmless polynomial factors in
$\ell$. Since $b\rho>1$, the resulting series converges geometrically.
The same Cauchy argument on a slightly larger compact set handles
$\boldsymbol u$ derivatives. This proves every interchange in the
initial half-space and then proves \eqref{sp:continuation} by continuation.
For $|U|<1/2$, the only Hurwitz poles occur when $\ell=k+1$.
Subtracting exactly their $1/U$ terms gives \eqref{sp:germ} and
\eqref{sp:H}.
\end{proof}

The germ need not be jointly holomorphic at zero: $\cR_P$ need not
be divisible by $U$. Its restriction to every ray
$\boldsymbol u=s\boldsymbol w$ with $W:=\sum_jw_j\ne0$ is nevertheless
holomorphic at $s=0$, since $\cR_P(0)=0$. The auxiliary $b$ can change
$\cR_P$ off $U=0$ by a multiple of $U$, absorbed into $\cH_P$.
Its restriction to the polar hyperplane, and the tangent coefficient
identities below, are intrinsic.

\subsection{Polynomiality in normalized exponents and exact finite differences}

For $\boldsymbol\alpha$ satisfying $\sum_j\alpha_j=1$, set
\begin{equation}\label{sp:F}
 \mathfrak F_r(\boldsymbol\alpha)
   =\left.\frac{d^r}{ds^r}
          \cZ_P(s\boldsymbol\alpha;\boldsymbol a)\right|_{s=0},
 \qquad r\ge0.
\end{equation}
Complex normalized exponents are allowed by continuation. Write
$\cH_{P,[q]}$ and $\cR_{P,[q]}$ for the homogeneous Taylor parts.

\begin{theorem}[Every normalized spectral jet is a polynomial]
\label{sp:thm:polynomial}
One has
\begin{align}
 \mathfrak F_r(\boldsymbol\alpha)
  &=r!\{\cH_{P,[r]}(\boldsymbol\alpha)
                +\cR_{P,[r+1]}(\boldsymbol\alpha)\},\label{sp:homogeneous}\\
 \cR_{P,[r+1]}(\boldsymbol\alpha)
  &=\frac{(-1)^{r+1}}{(r+1)!}
      \sum_{k=0}^D p_k[z^{k+1}]L_{\boldsymbol\alpha}(z)^{r+1}.
 \label{sp:Rhomogeneous}
\end{align}
Thus $\mathfrak F_r$ is a polynomial of degree at most $r+1$ on the
affine hyperplane $\sum\alpha_j=1$. Its component of degree $r+1$
is entirely local and vanishes when $r>D$.
\end{theorem}
\begin{proof}
Substitute $U=s$ in \eqref{sp:germ} and extract $s^r$.
Homogeneous extraction in $e^{-L_{\boldsymbol\alpha}}$ gives
\eqref{sp:Rhomogeneous}. The degree assertions follow immediately.
\end{proof}

For a function on the normalized hyperplane, put
$\Delta_{\boldsymbol v}f(\boldsymbol\alpha)
=f(\boldsymbol\alpha+\boldsymbol v)-f(\boldsymbol\alpha)$ when
$\sum_jv_j=0$. The notation $[n^{-1}]f(n)$ means the coefficient
of $n^{-1}$ in the Laurent expansion at large $n$; it is the negative
of the complex-analytic residue at infinity when that residue applies.

\begin{theorem}[Local polarization identity]\label{sp:thm:polarization}
Let $\boldsymbol v_1,\ldots,\boldsymbol v_{r+1}$ be tangent vectors
with $\sum_jv_{ij}=0$, and put
$V_i(n)=\sum_jv_{ij}\log(n+a_j)$. Then
\begin{align}\label{sp:finite-difference}
 &\Delta_{\boldsymbol v_1}\cdots\Delta_{\boldsymbol v_{r+1}}
       \mathfrak F_r(\boldsymbol\alpha)\notag\\
 &\quad=D_{\boldsymbol v_1}\cdots D_{\boldsymbol v_{r+1}}
       \mathfrak F_r(\boldsymbol\alpha)
       =(-1)^{r+1}r!\,[n^{-1}]P(n)\prod_{i=1}^{r+1}V_i(n).
\end{align}
All mixed finite differences of order $r+2$ vanish. If $\deg P=r$
and $p$ is its leading coefficient, the right side reduces to
\begin{equation}\label{sp:leading}
 (-1)^{r+1}r!p
             \prod_{i=1}^{r+1}\left(\sum_jv_{ij}a_j\right).
\end{equation}
\end{theorem}

\begin{proof}
The $r+1$ tangent differentiations kill $\cH_{P,[r]}$ and polarize
$L_{\boldsymbol\alpha}^{r+1}$, supplying a factor $(r+1)!$ in
\eqref{sp:Rhomogeneous}. Since $\sum_jv_{ij}=0$,
\[
 L_{\boldsymbol v_i}((n+b)^{-1})
   =\sum_jv_{ij}\log(n+a_j)=V_i(n)=O(n^{-1}).
\]
The coefficient $\sum_k p_k[z^{k+1}]\prod_iL_{\boldsymbol v_i}(z)$
is therefore $[n^{-1}]P(n)\prod_i V_i(n)$. Translation from $n+b$
to $n$ preserves this coefficient. For a polynomial of degree at most
$r+1$, its mixed difference of order $r+1$ equals its corresponding
constant mixed derivative; for example integrate that derivative over
the unit $(r+1)$-cube. This proves \eqref{sp:finite-difference}.
Finally $V_i(n)=n^{-1}\sum_jv_{ij}a_j+O(n^{-2})$ proves
\eqref{sp:leading}.
\end{proof}

Normalization cannot be omitted. For arbitrary $\boldsymbol w$ with
$W\ne0$, define $Z_{\boldsymbol w}(s)=\cZ_P(s\boldsymbol w)$.
Then
\begin{equation}\label{sp:scaling}
 Z_{\boldsymbol w}^{(r)}(0)
       =W^r\mathfrak F_r(\boldsymbol w/W).
\end{equation}
The degree theorem concerns factor exponents; it is not a polynomiality
assertion in the shifts $a_j$ themselves.

\subsection{A full Hurwitz--Stieltjes formula for each jet}

Put $A_{h,\ell}=[z^\ell]L_{\boldsymbol\alpha}(z)^h$, with
$A_{0,0}=1$ and $A_{0,\ell}=0$ for $\ell>0$.

\begin{corollary}[Normally convergent coefficient formula]\label{sp:cor:alljets}
For every $r\ge0$,
\begin{align}\label{sp:alljets}
 \mathfrak F_r(\boldsymbol\alpha)
 =\sum_{k=0}^Dp_k\Bigg\{&\zeta^{(r)}(-k,b)
 +\sum_{\substack{\ell\ge1\\\ell\ne k+1}}\sum_{h=1}^r
   (-1)^h\binom rh A_{h,\ell}\zeta^{(r-h)}(\ell-k,b)\notag\\
 &+\frac{(-1)^{r+1}}{r+1}A_{r+1,k+1}
   +(-1)^r\sum_{h=1}^r\binom rh A_{h,k+1}\gamma_{r-h}(b)
                                      \Bigg\}.
\end{align}
Empty sums vanish. The generalized Stieltjes constants occur linearly,
with indices at most $r-1$.
\end{corollary}
\begin{proof}
Multiply the expansion of $e^{-sL_{\boldsymbol\alpha}(z)}$ by the
Taylor expansion of each nonresonant Hurwitz factor in
\eqref{sp:continuation}. For $\ell=k+1$, multiply it by
$s^{-1}+\sum_q(-1)^q\gamma_q(b)s^q/q!$. The coefficient of $s^r$
receives a contribution from the exponential coefficient of order
$r+1$, which yields the third term in \eqref{sp:alljets}. Normal
convergence was proved in Theorem~\ref{sp:thm:germ}.
\end{proof}

This additional order-$r+1$ coefficient must be retained before a
spectral pole is canceled. Generalized harmonic numbers provide an
alternative finite description of all binomial coefficients required:
\begin{equation}\label{sp:harmonic}
 \frac{(x)_q}{q!}
 =\frac xq\exp\left(\sum_{j\ge1}
                   \frac{(-1)^{j+1}H_{q-1}^{(j)}}j x^j\right),
 \qquad q\ge1.
\end{equation}
This follows from $(x)_q/q!=(x/q)\prod_{j=1}^{q-1}(1+x/j)$.
Use it in \eqref{sp:Cpoch} and truncate at the desired spectral order.
Thus both the Stieltjes and generalized-harmonic content are explicit.

\subsection{First derivative: the classical weighted anomaly and Gamma factors}

Set
\begin{equation}\label{sp:single}
 \zeta_P(s,a)=\sum_{n\ge0}\frac{P(n)}{(n+a)^s}
       =\sum_{k=0}^D\frac{P^{(k)}(-a)}{k!}\zeta(s-k,a),
\end{equation}
where continuation is understood, and define
$\delta_{\boldsymbol w}=\sum_jw_j\zeta_P'(0,a_j)
-Z_{\boldsymbol w}'(0)$. The first-derivative instance of the calculus
recovers
\begin{equation}\label{sp:anomaly}
 \boxed{\delta_{\boldsymbol w}
   =\frac1{2W}\sum_{i<j}w_iw_j[n^{-1}]P(n)
                 \log^2\frac{n+a_i}{n+a_j}.}
\end{equation}
Indeed the nonlocal terms of $\mathfrak F_1$ are linear in the
normalized weights and cancel against their interpolation at the
vertices $\boldsymbol e_j$. The local difference is
\[
 \frac12\sum_kp_k[z^{k+1}]
 \left\{\sum_j\alpha_jL_j^2-
           \left(\sum_j\alpha_jL_j\right)^2\right\}
 =\frac12\sum_{i<j}\alpha_i\alpha_j
              \sum_kp_k[z^{k+1}](L_i-L_j)^2.
\]
Multiplication by $W$ gives \eqref{sp:anomaly}.

The sign convention is fixed by the determinant identity
\begin{equation}\label{sp:determinant}
 e^{-Z_{\boldsymbol w}'(0)}
   =e^{\delta_{\boldsymbol w}}
                     \prod_j e^{-w_j\zeta_P'(0,a_j)}.
\end{equation}
For $P(n)=\sum_{d=0}^Dq_dn^d$, the finite coefficient in
\eqref{sp:anomaly} is
\begin{equation}\label{sp:pairpolynomial}
 \sum_{d=1}^Dq_d(-1)^{d+1}
   \sum_{\ell=1}^d
   \frac{(a_i^\ell-a_j^\ell)
         (a_i^{d+1-\ell}-a_j^{d+1-\ell})}{\ell(d+1-\ell)}.
\end{equation}
For $P=1$ it is zero; for $P=n,n^2,n^3$ it is, respectively,
\begin{align*}
 &(a_i-a_j)^2,\qquad -(a_i-a_j)^2(a_i+a_j),\\
 &\frac{(a_i-a_j)^2(11a_i^2+14a_ia_j+11a_j^2)}{12}.
\end{align*}
For shifts $(1,2,3)$ and weights $(1,2,1)$, this gives
$\delta=1$ for $P=n$ and $\delta=-4$ for $P=n^2$.

With $G$ denoting the Barnes $G$ function normalized by $G(1)=1$,
$G(a+1)=\Gamma(a)G(a)$, the Hurwitz--Barnes identity gives, when $P(n)=n$,
\begin{equation}\label{sp:Barnes}
 \zeta_P'(0,a)=\zeta'(-1)-\log G(a+1)+\frac a2\log(2\pi).
\end{equation}
For a direct normalization check, differentiate both sides using
Hurwitz shift differentiation and the Barnes integral formula
\cite[Eq.~5.17.4]{dlmf}. Both derivatives are
$a-1/2-a\psi(a)$, and their values at $a=1$ agree. This proves the
displayed relation on $(0,\infty)$.
For two unit-weight factors the complete determinant is therefore
\begin{equation}\label{sp:Barnesdet}
 e^{-Z_{(1,1)}'(0)}
 =\frac{e^{(a-b)^2/4-2\zeta'(-1)}G(a+1)G(b+1)}
             {(2\pi)^{(a+b)/2}}.
\end{equation}
Here $a,b>0$ are the two shifts; $b$ in this last formula need not be
the expansion base used earlier. All logarithms of $G$ are real on
these positive arguments.

\subsection{A six-term identity for second derivatives}

\begin{theorem}[Quadratic-weight Vandermonde evaluation]
\label{sp:thm:sixterm}
Let $P$ be a monic quadratic polynomial and, for $a,b>0$, define
\begin{equation}\label{sp:Tdef}
 \mathfrak T_P(a,b)=\left.\frac{d^2}{ds^2}
     \sum_{n\ge0}P(n)\bigl((n+a)^2(n+b)\bigr)^{-s}
                                      \right|_{s=0},
\end{equation}
by meromorphic continuation of the initially convergent series. Then
for every $a,b,c>0$,
\begin{align}\label{sp:sixterm}
 &\mathfrak T_P(a,b)+\mathfrak T_P(b,c)+\mathfrak T_P(c,a)
 -\mathfrak T_P(b,a)-\mathfrak T_P(c,b)-\mathfrak T_P(a,c)\notag\\
 &\hspace{30mm}=\frac23(a-b)(b-c)(c-a).
\end{align}
For a quadratic of leading coefficient $p$, multiply the right side
by $p$. For degree at most one, the right side is zero.
\end{theorem}

\begin{proof}
Use the three shifts $(a,b,c)$, take
$\boldsymbol\alpha_0=(1,1,1)/3$, and take tangent vectors
\[
 \boldsymbol v_1=(1,-1,0)/3,\quad
 \boldsymbol v_2=(0,1,-1)/3,\quad
 \boldsymbol v_3=(-1,0,1)/3.
\]
Their sum is zero, so the empty and full subsets cancel in the
third finite difference of $\mathfrak F_2$. Equation~\eqref{sp:leading}
gives
\[
 \sum_{i=1}^3\{\mathfrak F_2(\boldsymbol\alpha_0+\boldsymbol v_i)
      -\mathfrak F_2(\boldsymbol\alpha_0-\boldsymbol v_i)\}
 =-\frac2{27}(a-b)(b-c)(c-a).
\]
Each argument is a permutation of $(2/3,1/3,0)$. By
\eqref{sp:scaling} its value is the corresponding $\mathfrak T_P/9$.
Reading the six arguments gives the negative of the left side of
\eqref{sp:sixterm}, divided by nine. This proves the formula and its
sign. The lower-degree assertions follow from the coefficient
$[n^{-1}]$ in \eqref{sp:finite-difference}.
\end{proof}

For $(a,b,c)=(1,2,3)$ the value is $4/3$, independent of the lower
two coefficients of $P$. This is a cancellation of all nonlocal
second-derivative constants; applying only a first-derivative anomaly
formula would not justify it.

\subsection{An alternating Vandermonde identity with any number of factors}

The six-term evaluation is the first nontrivial higher-derivative member
of a general alternating family. Define
\[
 N=\binom m2,\qquad
 \Delta(\boldsymbol x)=\prod_{1\le i<j\le m}(x_j-x_i),\qquad
 K_m=\prod_{k=0}^{m-1}k!,
\]
and let $(\sigma\boldsymbol w)_j=w_{\sigma(j)}$.

\begin{theorem}[General spectral alternant]\label{sp:thm:alternant}
For $m\ge2$, shifts $a_j>0$, and weights with $W=\sum_jw_j\ne0$,
\begin{align}\label{sp:alternant}
 &\sum_{\sigma\in S_m}\operatorname{sgn}(\sigma)
       Z_{\sigma\boldsymbol w}^{(N-1)}(0)\notag\\
 &\quad=\frac{(-1)^N(N-1)!}{WK_m}\Delta(\boldsymbol w)
     [n^{-1}]P(n)\prod_{i<j}\log\frac{n+a_j}{n+a_i}.
\end{align}
If $\deg P=N-1$ with leading coefficient $p$, this is the complete
algebraic evaluation
\begin{equation}\label{sp:alternant-leading}
 \boxed{\sum_{\sigma\in S_m}\operatorname{sgn}(\sigma)
       Z_{\sigma\boldsymbol w}^{(N-1)}(0)
 =\frac{(-1)^N(N-1)!p}{WK_m}
                 \Delta(\boldsymbol w)\Delta(\boldsymbol a).}
\end{equation}
For degree less than $N-1$, the alternating sum is zero.
\end{theorem}

\begin{proof}
Every alternating polynomial in $m$ variables is divisible by their
Vandermonde polynomial, of degree $N$. Therefore alternation kills
the degree-at-most-$N-1$ term $\cH_{P,[N-1]}$ in
\eqref{sp:homogeneous}. The elementary alternant identity is
\begin{equation}\label{sp:elementary-alternant}
 \sum_{\sigma\in S_m}\operatorname{sgn}(\sigma)
       \left(\sum_j\alpha_{\sigma(j)}y_j\right)^N
   =\frac{N!}{K_m}\Delta(\boldsymbol\alpha)\Delta(\boldsymbol y).
\end{equation}
To prove it, expand by the multinomial theorem. A monomial with
repeated exponents disappears under alternation. The only distinct
nonnegative exponent pattern of total degree $N$ is a permutation
of $0,1,\ldots,m-1$. The surviving determinant in the $\alpha_j$
and the corresponding determinant in the $y_j$ give the two
Vandermondes, with the positive orientation shown in
\eqref{sp:elementary-alternant} and multinomial factor $N!/K_m$.

Apply this identity to $y_j=L_j(z)$ in
\eqref{sp:Rhomogeneous}. The normalized alternating jet is
\[
 \frac{(-1)^N(N-1)!}{K_m}\Delta(\boldsymbol\alpha)
           \sum_kp_k[z^{k+1}]\prod_{i<j}(L_j(z)-L_i(z)).
\]
At $z=(n+b)^{-1}$, every difference is
$\log((n+a_j)/(n+a_i))$. The coefficient is therefore the one
in \eqref{sp:alternant}. Use \eqref{sp:scaling} and
$\Delta(\boldsymbol w/W)=W^{-N}\Delta(\boldsymbol w)$ to obtain
the unnormalized factor $1/W$. Finally, the logarithmic product
starts with $\Delta(\boldsymbol a)n^{-N}$, proving
\eqref{sp:alternant-leading} and the lower-degree vanishing.
\end{proof}

For $m=3$ and $\boldsymbol w=(2,1,0)$, equation
\eqref{sp:alternant-leading} is exactly \eqref{sp:sixterm}.
For $m=4$, $N=6$ and $(N-1)!/K_m=10$: a signed sum of
twenty-four fifth spectral derivatives equals
$10p\Delta(\boldsymbol w)\Delta(\boldsymbol a)/W$.
For example, with $\boldsymbol w=(3,2,1,0)$,
$\boldsymbol a=(1,2,3,4)$ and any monic degree-five $P$, that
sum is exactly $240$. No individual fifth derivative is evaluated
to obtain this algebraic cancellation.

\subsection{Shift derivatives, polygammas, and normalized primitives}

The single-factor spectral derivative has the exact shift derivative
\begin{equation}\label{sp:shiftderivative}
 \frac{d}{da}\zeta_P'(0,a)
  =-\FP_{s=1}\zeta_P(s,a)
  =P(-a)\psi(a)
      -\sum_{k=1}^D\frac{P^{(k)}(-a)}{k!}\zeta(1-k,a).
\end{equation}
To prove it, differentiate the convergent Dirichlet series to obtain
$\partial_a\zeta_P(s,a)=-s\zeta_P(s+1,a)$. The latter zeta factor
has residue $P(-a)$ at $s=0$. Multiplication by $s$ before
differentiation shows that its finite part, rather than an undefined
value at the pole, is required. Since
$\zeta(1-k,a)=-B_k(a)/k$, the result is a polynomial times digamma
plus a polynomial. Repeated differentiation gives explicit polygamma
identities; integration over $[a_0,a_1]\subset(0,\infty)$ gives the
normalized primitive $\zeta_P'(0,a_1)-\zeta_P'(0,a_0)$.

Combining \eqref{sp:shiftderivative} with the finite anomaly polynomial
gives all shift derivatives of the product determinant. For example,
for $P=n$ and two unit weights,
\[
 \partial_a Z_{(1,1)}'(0)=\frac{a+b-1}2-a\psi(a).
\]

\subsection{Logarithmic multiplicities and higher-order poles}

Fix the same expansion base $b$ as in \eqref{sp:base}, and for an
integer $h\ge0$ consider the distinct spectrum
\begin{equation}\label{sp:logdef}
 \cZ_{P,h}(\boldsymbol u;\boldsymbol a;b)
  =\sum_{n\ge0}P(n)\log^h(n+b)\prod_j(n+a_j)^{-u_j}.
\end{equation}
Changing this $b$ changes the logarithmic multiplicity itself.
It must not be treated as a freely variable auxiliary base in this
definition.

\begin{theorem}[Higher-pole finite parts and their local identities]
\label{sp:thm:log}
Near zero,
\begin{equation}\label{sp:loggerm}
 \cZ_{P,h}(\boldsymbol u;\boldsymbol a;b)
   =\frac{h!\cR_P(\boldsymbol u)}{U^{h+1}}
                                +\cH_{P,h}(\boldsymbol u),
\end{equation}
where $\cH_{P,h}$ is holomorphic. For normalized $\boldsymbol\alpha$,
define the Laurent jet
\begin{equation}\label{sp:logjet}
 \mathfrak J_{r,h}(\boldsymbol\alpha)
  :=r![s^r]\cZ_{P,h}(s\boldsymbol\alpha;\boldsymbol a;b),
 \qquad r\ge0.
\end{equation}
It is a polynomial of degree at most $q:=r+h+1$, and
\begin{align}
 \mathfrak J_{r,h}
   &=r!\{\cH_{P,h,[r]}+h!\cR_{P,[r+h+1]}\},\label{sp:loghomogeneous}\\
 \Delta_{\boldsymbol v_1}\cdots\Delta_{\boldsymbol v_q}
       \mathfrak J_{r,h}(\boldsymbol\alpha)
   &=h!(-1)^q r!\,[n^{-1}]P(n)\prod_{i=1}^q V_i(n)
 \label{sp:logdifference}
\end{align}
for tangent vectors as before. All differences of order $q+1$ vanish.
\end{theorem}

\begin{proof}
The proof of \eqref{sp:continuation} is unchanged, except that each
Hurwitz factor is replaced by $(-1)^h\zeta^{(h)}$. At resonance,
\[
 (-1)^h\zeta^{(h)}(1+U,b)
  =\frac{h!}{U^{h+1}}
          +\sum_{j\ge0}\frac{(-1)^j\gamma_{j+h}(b)}{j!}U^j.
\]
Subtract the displayed pole to prove \eqref{sp:loggerm}. Normal
convergence persists after these fixed derivatives. Taking $[s^r]$
gives \eqref{sp:loghomogeneous}, and polarizing its degree-$q$ part
as in Theorem~\ref{sp:thm:polarization} proves
\eqref{sp:logdifference}.
\end{proof}

If $r+h>D$, the residue contribution in \eqref{sp:loghomogeneous}
vanishes, sharpening the degree bound to $r$. In particular, if $h>D$,
the constant finite part is independent of the factor shifts and
normalized weights and equals
$(-1)^h\sum_kp_k\zeta^{(h)}(-k,b)$, even though negative Laurent
powers may remain.

There is also a higher-pole version of Theorem~\ref{sp:thm:alternant}.
Let $N=\binom m2$, $0\le h\le N-1$, and $r=N-h-1$. The identical
alternant proof applied to \eqref{sp:loghomogeneous} gives
\begin{align}\label{sp:logalternant}
 &\sum_{\sigma\in S_m}\operatorname{sgn}(\sigma)\,
      r![s^r]\cZ_{P,h}(s\sigma\boldsymbol w;\boldsymbol a;b)\notag\\
 &\quad=\frac{h!(-1)^N r!}{W^{h+1}K_m}\Delta(\boldsymbol w)
       [n^{-1}]P(n)\prod_{i<j}\log\frac{n+a_j}{n+a_i}.
\end{align}
For $\deg P=N-1$ the final coefficient is
$p\Delta(\boldsymbol a)$. Thus the alternating algebraic evaluation
persists for logarithmic multiplicities, with the pole order determining
the required Laurent jet and the power of $W$.

For $h>0$ the ray generally has a pole of order at most $h$.
Consequently \eqref{sp:logjet} is a Laurent coefficient, not an
ordinary derivative of the unmodified spectrum at zero. With
$\mu=\sum_j\alpha_j(a_j-b)$, two concrete constant terms are
\begin{align}
 \FP_{s=0}\sum_{n\ge0}n\log(n+b)\prod_j(n+a_j)^{-\alpha_js}
  &=-\zeta'(-1,b)+b\zeta'(0,b)+\frac{\mu^2}2,
 \label{sp:logexample1}\\
 \FP_{s=0}\sum_{n\ge0}n^2\log^2(n+b)\prod_j(n+a_j)^{-\alpha_js}
  &=\zeta''(-2,b)-2b\zeta''(-1,b)+b^2\zeta''(0,b)
       -\frac{\mu^3}3.
 \label{sp:logexample2}
\end{align}
The positive-degree dependence of these finite parts is completely
algebraic despite the higher spectral poles. These formulas resolve
the concrete logarithmic-multiplicity extension of the local calculus;
arbitrary nonpolynomial spectra require a separate continuation theorem.

